\section{روش دالامبر و معادلۀ انتقال حرارت یک‌بعدی}
\def\EMfootchapter{فصل ۵ \ \lr{\textbullet}\ روش دالامبر و معادلۀ انتقال حرارت یک‌بعدی}
\def\EMtitlesize{26}\def\EMtitleskip{31}
\begin{frame}{فصل ۵: روش دالامبر و معادلۀ انتقال حرارت یک‌بعدی}
{\fontsize{20}{26}\selectfont\bfseries\color{SLdef}\EMalign مطالب این فصل\par}
\vspace{10bp}
\EMitem{\lr{\textbullet}}{\dimexpr\linewidth-26bp\relax}{روش دالامبر در حل معادلۀ موجی یک‌بعدی}
\EMitem{\lr{\textbullet}}{\dimexpr\linewidth-26bp\relax}{معادلۀ دیفرانسیل انتقال حرارت}
\EMitem{\lr{\textbullet}}{\dimexpr\linewidth-26bp\relax}{حل معادلۀ دیفرانسیل انتقال حرارت یک‌بعدی}
\EMitem{\lr{\textbullet}}{\dimexpr\linewidth-26bp\relax}{معادلۀ انتقال حرارت یک‌بعدی بر روی مفتول نامتناهی}
\EMitem{\lr{\textbullet}}{\dimexpr\linewidth-26bp\relax}{تابع خطا}
\end{frame}

\def\EMtitlesize{26}\def\EMtitleskip{31}
\begin{frame}{روش دالامبر در حل معادلۀ موجی یک‌بعدی}
\begin{uncoverenv}<2->
\EMeq{\begin{cases}\dfrac{\partial^2u}{\partial t^2}=c^2\dfrac{\partial^2u}{\partial x^2}\\\noalign{\vskip 2mm} u(0,t)=u(L,t)=0\\ u(x,0)=f(x)\\ u_t(x,0)=g(x)\end{cases}}
\end{uncoverenv}
\EMgap
\begin{uncoverenv}<3->
\EMp{با استفاده از تغییر متغیرها به شکل زیر معادلۀ دیفرانسیل را بازنویسی و حل می‌کنیم.}
\vspace{3bp}
\EMeq{u(x,t) \Rightarrow u(\varphi,\psi),\qquad \begin{cases}\varphi=x+ct\\ \psi=x-ct\end{cases}}
\end{uncoverenv}
\EMgap
\begin{uncoverenv}<4->
\EMeq{\frac{\partial u}{\partial x} =\frac{\partial u}{\partial\varphi}\frac{\partial\varphi}{\partial x}+\frac{\partial u}{\partial\psi}\frac{\partial\psi}{\partial x} =\frac{\partial u}{\partial\varphi}+\frac{\partial u}{\partial\psi}}
\end{uncoverenv}
\EMgap
\begin{uncoverenv}<5->
\EMeq{\frac{\partial u}{\partial t} =\frac{\partial u}{\partial\varphi}\frac{\partial\varphi}{\partial t}+\frac{\partial u}{\partial\psi}\frac{\partial\psi}{\partial t} =c\frac{\partial u}{\partial\varphi}-c\frac{\partial u}{\partial\psi}}
\end{uncoverenv}
\end{frame}

\def\EMtitlesize{26}\def\EMtitleskip{31}
\begin{frame}{روش دالامبر در حل معادلۀ موجی یک‌بعدی (ادامه)}
\begin{uncoverenv}<2->
\EMeq{\frac{\partial^2u}{\partial x^2} =\frac\partial{\partial x}\left(\frac{\partial u}{\partial\varphi}+\frac{\partial u}{\partial\psi}\right) =\left(\frac{\partial^2u}{\partial\varphi^2}+\frac{\partial^2u}{\partial\varphi\partial\psi}\right)\cdot1+\left(\frac{\partial^2u}{\partial\psi\partial\varphi}+\frac{\partial^2u}{\partial\psi^2}\right)\cdot1 =\frac{\partial^2u}{\partial\varphi^2}+2\frac{\partial^2u}{\partial\varphi\partial\psi}+\frac{\partial^2u}{\partial\psi^2}}
\end{uncoverenv}
\EMgap
\begin{uncoverenv}<3->
\EMp{(با استفاده از \EMim{\dfrac{\partial^2u}{\partial\varphi\partial\psi}=\dfrac{\partial^2u}{\partial\psi\partial\varphi}}.)}
\end{uncoverenv}
\end{frame}

\def\EMtitlesize{26}\def\EMtitleskip{31}
\begin{frame}{روش دالامبر در حل معادلۀ موجی یک‌بعدی (ادامه)}
\begin{EMbox}{exe}{1}{}
\EMp{به طریق مشابه تساوی زیر را ثابت کنید:}
\vspace{3bp}
\EMeq{\frac{\partial^2u}{\partial t^2} =c^2\left(\frac{\partial^2u}{\partial\varphi^2}-2\frac{\partial^2u}{\partial\varphi\partial\psi}+\frac{\partial^2u}{\partial\psi^2}\right)}
\end{EMbox}
\begin{uncoverenv}<2->
\EMp{با قرار دادن دو رابطۀ بالا در معادلۀ دیفرانسیل موجی یک‌بعدی:}
\vspace{3bp}
\EMeq{c^2\left(\frac{\partial^2u}{\partial\varphi^2}-2\frac{\partial^2u}{\partial\varphi\partial\psi}+\frac{\partial^2u}{\partial\psi^2}\right) =c^2\left(\frac{\partial^2u}{\partial\varphi^2}+2\frac{\partial^2u}{\partial\varphi\partial\psi}+\frac{\partial^2u}{\partial\psi^2}\right)\quad \Longrightarrow\quad \frac{\partial^2u}{\partial\varphi\partial\psi} =0}
\end{uncoverenv}
\EMgap
\begin{uncoverenv}<3->
\EMp{\textbf{حل عمومی:}}
\vspace{3bp}
\EMeq{\frac{\partial u}{\partial\varphi} =\phi_1(\varphi)\  \Rightarrow\ u(\varphi,\psi) =\int\phi_1(\varphi)\,d\varphi+\Psi(\psi) =\Phi(\varphi)+\Psi(\psi)\quad \Longrightarrow\quad u(x,t) \\[4pt] =\Phi(x+ct)+\Psi(x-ct)}
\end{uncoverenv}
\EMgap
\begin{uncoverenv}<4->
\EMp{\textbf{اعمال شرایط اولیه:}}
\vspace{3bp}
\EMeq{u(x,0) =\Phi(x)+\Psi(x) =f(x)}
\end{uncoverenv}
\end{frame}

\def\EMtitlesize{26}\def\EMtitleskip{31}
\begin{frame}{روش دالامبر در حل معادلۀ موجی یک‌بعدی (ادامه)}
\begin{uncoverenv}<2->
\EMeq{u_t(x,t) =\frac{\partial\Phi}{\partial\varphi}\frac{\partial\varphi}{\partial t}+\frac{\partial\Psi}{\partial\psi}\frac{\partial\psi}{\partial t} =c\Phi'(\varphi)-c\Psi'(\psi) =c\Phi'(x+ct)-c\Psi'(x-ct)}
\end{uncoverenv}
\EMgap
\begin{uncoverenv}<3->
\EMeq{u_t(x,0) =c\Phi'(x)-c\Psi'(x) =c\big[\Phi'(x)-\Psi'(x)\big] =g(x)}
\end{uncoverenv}
\EMgap
\begin{uncoverenv}<4->
\EMeq{\begin{cases}\Phi(x)+\Psi(x)=f(x)\\ \Phi'(x)-\Psi'(x)=\dfrac1cg(x)\end{cases} \Rightarrow\begin{cases}\Phi'(x)+\Psi'(x)=f'(x)\\ \Phi'(x)-\Psi'(x)=\dfrac1cg(x)\end{cases} \Rightarrow\begin{cases}\Phi'(x)=\dfrac12f'(x)+\dfrac1{2c}g(x)\\\noalign{\vskip 2mm} \Psi'(x)=\dfrac12f'(x)-\dfrac1{2c}g(x)\end{cases}}
\end{uncoverenv}
\EMgap
\begin{uncoverenv}<5->
\EMeq{\begin{cases}\Phi(x)=\dfrac12f(x)+\dfrac1{2c}\displaystyle\int_0^xg(s)\,ds+k_1\\\noalign{\vskip 2mm} \Psi(x)=\dfrac12f(x)-\dfrac1{2c}\displaystyle\int_0^xg(s)\,ds+k_2\end{cases}}
\end{uncoverenv}
\EMgap
\begin{uncoverenv}<6->
\EMeq{u(x,t) =\frac12\left(f(x+ct)+\frac1c\int_0^{x+ct}g(s)\,ds\right)+\frac12\left(f(x-ct)-\frac1c\int_0^{x-ct}g(s)\,ds\right)+k_1+k_2}
\end{uncoverenv}
\end{frame}

\def\EMtitlesize{26}\def\EMtitleskip{31}
\begin{frame}{روش دالامبر در حل معادلۀ موجی یک‌بعدی (ادامه)}
\begin{EMbox}{imp}{1}{}
\begin{uncoverenv}<2->
\EMeq{u(x,t) =\frac12\Big[f(x+ct)+f(x-ct)\Big]+\frac1{2c}\int_{x-ct}^{x+ct}g(s)\,ds+K,\qquad u(0,0) =0\  \Rightarrow\ K =0}
\end{uncoverenv}
\EMgap
\begin{uncoverenv}<3->
\EMeq{u(x,t) =\frac12\Big[f(x+ct)+f(x-ct)\Big]+\frac1{2c}\int_{x-ct}^{x+ct}g(s)\,ds}
\end{uncoverenv}
\end{EMbox}
\end{frame}

\def\EMtitlesize{26}\def\EMtitleskip{31}
\begin{frame}{روش دالامبر در حل معادلۀ موجی یک‌بعدی (ادامه)}
\begin{EMbox}{exe}{1}{}
\EMp{ثابت کنید که اگر بخواهیم شرایط مرزی برقرار باشند بایستی در رابطۀ بالا توابع \EMim{f(x)} و \EMim{g(x)} هر دو تناوبی با دورۀ تناوب \EMim{2L} بوده و فرد باشند.}
\end{EMbox}
\end{frame}

\def\EMtitlesize{26}\def\EMtitleskip{31}
\begin{frame}{معادلۀ دیفرانسیل انتقال حرارت}
\begin{uncoverenv}<2->
\EMeq{\frac{\partial u}{\partial t} =c^2\frac{\partial^2u}{\partial x^2},\qquad \frac{\partial u}{\partial t} =c^2\left(\frac{\partial^2u}{\partial x^2}+\frac{\partial^2u}{\partial y^2}\right),\qquad \frac{\partial u}{\partial t} =c^2\left(\frac{\partial^2u}{\partial x^2}+\frac{\partial^2u}{\partial y^2}+\frac{\partial^2u}{\partial z^2}\right)}
\end{uncoverenv}
\end{frame}

\def\EMtitlesize{26}\def\EMtitleskip{31}
\begin{frame}{معادلۀ دیفرانسیل انتقال حرارت یک‌بعدی}
\begin{minipage}[t]{500bp}\vspace{0pt}
\EMgap
\begin{uncoverenv}<3->
\EMp{برای حل یکتای این معادله نیاز به شرایط مرزی و شرایط اولیه داریم:}
\end{uncoverenv}
\EMgap
\begin{uncoverenv}<4->
\EMitem{\lr{\textbullet}}{\dimexpr\linewidth-26bp\relax}{شرایط مرزی (\lr{Boundary Conditions}): \EMim{u(0,t)=u(L,t)=0}}
\end{uncoverenv}
\EMgap
\begin{uncoverenv}<5->
\EMitem{\lr{\textbullet}}{\dimexpr\linewidth-26bp\relax}{شرایط اولیه (\lr{Initial Conditions}): \EMim{u(x,0)=f(x)}}
\end{uncoverenv}
\end{minipage}
\EMside{0.0bp}{\uncover<2->{\EMfigure{1.500}{m05-rod}{میلۀ یک‌بعدی \EMim{0\le x\le L} با توزیع دمای \EMim{u(x,t)}}}}
\end{frame}

\def\EMtitlesize{26}\def\EMtitleskip{31}
\begin{frame}{حل معادلۀ دیفرانسیل انتقال حرارت یک‌بعدی}
\begin{uncoverenv}<2->
\EMeq{\begin{cases}\dfrac{\partial u}{\partial t}=c^2\dfrac{\partial^2u}{\partial x^2}\\\noalign{\vskip 2mm} u(0,t)=u(L,t)=0\\ u(x,0)=f(x)\end{cases}}
\end{uncoverenv}
\EMgap
\begin{uncoverenv}<3->
\EMp{در اینجا برای حل این معادله از روش \textbf{جداسازی متغیرها} استفاده می‌کنیم. فرض: \EMim{u(x,t)=X(x)\times T(t)}.}
\vspace{3bp}
\EMeq{\frac{\partial^2u}{\partial x^2} =X''(x)T(t),\qquad \frac{\partial u}{\partial t} =X(x)T'(t)\quad \Longrightarrow\quad X(x)T'(t) =c^2X''(x)T(t)\quad \Longrightarrow\quad \frac{T'}{c^2T} =\frac{X''}X =k}
\end{uncoverenv}
\EMgap
\begin{uncoverenv}<4->
\EMp{\EMim{k} باید حتماً یک عدد منفی باشد (برای \EMim{k=0} و \EMim{k>0} دقیقاً مانند مسئلۀ موجی \EMim{X(x)=0} به دست می‌آید). با \EMim{k=-\lambda^2}:}
\vspace{3bp}
\EMeq{\frac{T'}{c^2T} =\frac{X''}X =-\lambda^2\quad \Longrightarrow\quad \begin{cases}X''+\lambda^2X=0\\ T'+c^2\lambda^2T=0\end{cases}}
\end{uncoverenv}
\EMgap
\begin{uncoverenv}<5->
\EMeq{X(x) =A\cos\lambda x+B\sin\lambda x,\quad X(0) =0 \Rightarrow A =0,\qquad X(L) =0 \Rightarrow B\sin\lambda L =0 \Rightarrow\lambda =\frac{n\pi}L =\lambda_n}
\end{uncoverenv}
\end{frame}

\def\EMtitlesize{26}\def\EMtitleskip{31}
\begin{frame}{حل معادلۀ دیفرانسیل انتقال حرارت یک‌بعدی (ادامه)}
\begin{uncoverenv}<2->
\EMeq{X_n(x) =B_n\sin\lambda_nx,\qquad T'+c^2\lambda_n^2T =0\  \Rightarrow\ T_n(t) =C_ne^{-(c\lambda_n)^2t}}
\end{uncoverenv}
\EMgap
\begin{uncoverenv}<3->
\EMeq{u_n(x,t) =X_n(x)\times T_n(t) =B_n\sin\!\left(\frac{n\pi}Lx\right)\times C_ne^{-\left(\frac{cn\pi}L\right)^2t}}
\end{uncoverenv}
\EMgap
\begin{uncoverenv}<4->
\EMeq{u(x,t) =\sum_{n=1}^{\infty}A_ne^{-(c\lambda_n)^2t}\sin(\lambda_nx),\qquad \lambda_n =\frac{n\pi}L}
\end{uncoverenv}
\EMgap
\begin{uncoverenv}<5->
\EMp{شرط اولیه \EMim{u(x,0)=f(x)}:}
\vspace{3bp}
\EMeq{u(x,0) =\sum_{n=1}^{\infty}A_n\sin\lambda_nx =f(x)}
\end{uncoverenv}
\end{frame}

\def\EMtitlesize{26}\def\EMtitleskip{31}
\begin{frame}{حل معادلۀ دیفرانسیل انتقال حرارت یک‌بعدی (ادامه)}
\begin{uncoverenv}<2->
\EMp{تابع \EMim{f(x)} با یک تابع تناوبی و فرد برابر شده است، بنابراین برای محاسبۀ ضرایب مجهول فوریه باید \EMim{f(x)} را یک تابع تناوبی (با دورۀ تناوب \EMim{2L}) و فرد فرض کنیم:}
\vspace{3bp}
\EMeq{A_n =\frac2{2L}\int_{-L}^{L}f(x)\sin\!\left(\frac{2\pi}{2L}nx\right)dx =\frac1L\int_{-L}^{L}f(x)\sin\!\left(\frac{n\pi}Lx\right)dx =\frac2L\int_0^Lf(x)\sin\!\left(\frac{n\pi}Lx\right)dx}
\end{uncoverenv}
\EMgap
\begin{EMbox}{imp}{2}{}
\begin{uncoverenv}<3->
\EMeq{u(x,t) =\sum_{n=1}^{\infty}A_ne^{-(c\lambda_n)^2t}\sin(\lambda_nx),\qquad \lambda_n =\frac{n\pi}L,\qquad A_n =\frac2L\int_0^Lf(x)\sin\!\left(\frac{n\pi}Lx\right)dx}
\end{uncoverenv}
\end{EMbox}
\end{frame}

\def\EMtitlesize{26}\def\EMtitleskip{31}
\begin{frame}{حل معادلۀ دیفرانسیل انتقال حرارت یک‌بعدی (ادامه)}
\begin{EMbox}{ex}{1}{}
\EMp{جواب معادلۀ انتقال حرارت یک‌بعدی زیر را به‌دست آورید.}
\vspace{3bp}
\EMeq{\frac{\partial u}{\partial t} =\frac{\partial^2u}{\partial x^2},\qquad u(0,t) =u(10,t) =0,\qquad u(x,0) =f(x) =\begin{cases}x&0\le x<5\\ 10-x&5\le x\le10\end{cases}}
\EMgap
\begin{uncoverenv}<2->
\EMeq{\begin{aligned}
A_n&=\frac2L\int_0^Lf(x)\sin\!\left(\frac{n\pi}Lx\right)dx=0.2\int_0^5x\sin\!\left(\frac{n\pi}{10}x\right)dx+0.2\int_5^{10}(10-x)\sin\!\left(\frac{n\pi}{10}x\right)dx\\
&=0.2\left[(x)\left(\frac{-10}{n\pi}\cos\frac{n\pi}{10}x\right)-(1)\left(\frac{-100}{n^2\pi^2}\sin\frac{n\pi}{10}x\right)\right]_0^5\\
&\quad+0.2\left[(10-x)\left(\frac{-10}{n\pi}\cos\frac{n\pi}{10}x\right)-(-1)\left(\frac{-100}{n^2\pi^2}\sin\frac{n\pi}{10}x\right)\right]_5^{10}=\frac{40}{n^2\pi^2}\sin\frac{n\pi}2
\end{aligned}}
\end{uncoverenv}
\EMgap
\begin{uncoverenv}<3->
\EMeq{u(x,t) =\sum_{n=1}^{\infty}\left(\frac{40}{n^2\pi^2}\sin\frac{n\pi}2\right)e^{-\left(\frac{cn\pi}{10}\right)^2t}\sin\!\left(\frac{n\pi}{10}x\right)}
\end{uncoverenv}
\end{EMbox}
\end{frame}

\def\EMtitlesize{26}\def\EMtitleskip{31}
\begin{frame}{حل معادلۀ دیفرانسیل انتقال حرارت یک‌بعدی (ادامه)}
\begin{EMbox}{ex}{1}{}
\begin{uncoverenv}<2->
\EMp{(با \EMim{c=1}.)}
\end{uncoverenv}
\end{EMbox}
\end{frame}

\def\EMtitlesize{26}\def\EMtitleskip{31}
\begin{frame}{حل معادلۀ دیفرانسیل انتقال حرارت یک‌بعدی (ادامه)}
\begin{EMbox}{ex}{1}{}
\EMp{جواب معادلۀ انتقال حرارت یک‌بعدی زیر را به‌دست آورید.}
\vspace{3bp}
\EMeq{\frac{\partial u}{\partial t} =c^2\frac{\partial^2u}{\partial x^2},\qquad u(0,t) =T_0,\qquad u(L,t) =T_1,\qquad u(x,0) =f(x)}
\EMgap
\begin{uncoverenv}<2->
\EMeq{u(x,t) =w(x,t)+v(x)\quad \Longrightarrow\quad \begin{cases}w(0,t)+v(0)=T_0\\ w(L,t)+v(L)=T_1\end{cases}}
\end{uncoverenv}
\EMgap
\begin{uncoverenv}<3->
\EMeq{\frac{\partial w}{\partial t} =c^2\left(\frac{\partial^2w}{\partial x^2}+\frac{\partial^2v}{\partial x^2}\right)}
\end{uncoverenv}
\end{EMbox}
\end{frame}

\def\EMtitlesize{26}\def\EMtitleskip{31}
\begin{frame}{حل معادلۀ دیفرانسیل انتقال حرارت یک‌بعدی (ادامه)}
\begin{EMbox}{ex}{1}{}
\begin{uncoverenv}<2->
\EMeq{\begin{cases}\dfrac{\partial^2v}{\partial x^2}=0\\ v(0)=T_0\\ v(L)=T_1\end{cases} \Rightarrow\ v(x) =\left(\frac{T_1-T_0}L\right)x+T_0,\qquad \begin{cases}\dfrac{\partial w}{\partial t}=c^2\dfrac{\partial^2w}{\partial x^2}\\ w(0,t)=w(L,t)=0\end{cases} \Rightarrow\ w(x,t) \\[4pt] =\sum_{n=1}^{\infty}A_ne^{-(c\lambda_n)^2t}\sin(\lambda_nx)}
\end{uncoverenv}
\EMgap
\begin{uncoverenv}<3->
\EMeq{u(x,t) =\underbrace{\sum_{n=1}^{\infty}A_ne^{-(c\lambda_n)^2t}\sin(\lambda_nx)}_{w(x,t)}+\underbrace{\left(\frac{T_1-T_0}L\right)x+T_0}_{v(x)},\qquad \lambda_n =\frac{n\pi}L}
\end{uncoverenv}
\end{EMbox}
\end{frame}

\def\EMtitlesize{26}\def\EMtitleskip{31}
\begin{frame}{حل معادلۀ دیفرانسیل انتقال حرارت یک‌بعدی (ادامه)}
\begin{EMbox}{ex}{1}{}
\begin{uncoverenv}<2->
\EMp{شرط اولیه:}
\vspace{3bp}
\EMeq{\sum_{n=1}^{\infty}A_n\sin\frac{n\pi}Lx =f(x)-\left(\frac{T_1-T_0}L\right)x-T_0\triangleq f_1(x)\  \Longrightarrow\ A_n \\[4pt] =\frac2L\int_0^L\left[f(x)-\left(\frac{T_1-T_0}L\right)x-T_0\right]\sin\frac{n\pi}Lx\,dx}
\end{uncoverenv}
\end{EMbox}
\end{frame}

\def\EMtitlesize{26}\def\EMtitleskip{31}
\begin{frame}{حل معادلۀ دیفرانسیل انتقال حرارت یک‌بعدی (ادامه)}
\begin{EMbox}{ex}{1}{}
\EMp{جواب معادلۀ انتقال حرارت یک‌بعدی زیر را به‌دست آورید.}
\vspace{3bp}
\EMeq{\frac{\partial u}{\partial t} =c^2\frac{\partial^2u}{\partial x^2},\qquad u_x(0,t) =u_x(L,t) =0,\qquad u(x,0) =f(x)}
\EMgap
\begin{uncoverenv}<2->
\EMp{برای حل این مسئله از روش \textbf{جداسازی متغیرها} استفاده می‌کنیم: \EMim{u(x,t)=X(x)\times T(t)}.}
\vspace{3bp}
\EMeq{\frac{T'}{c^2T} =\frac{X''}X =-\lambda^2\  \Rightarrow\ \begin{cases}X''+\lambda^2X=0\ \Rightarrow\ X(x)=A\cos\lambda x+B\sin\lambda x\\ T'+c^2\lambda^2T=0\end{cases}}
\end{uncoverenv}
\EMgap
\begin{uncoverenv}<3->
\EMeq{X'(x) =-A\lambda\sin\lambda x+B\lambda\cos\lambda x,\qquad X'(0) =0 \Rightarrow B =0,\qquad X'(L) =0 \Rightarrow-A\lambda\sin\lambda L =0 \Rightarrow\lambda L \\[4pt] =n\pi \Rightarrow\lambda_n =\frac{n\pi}L\quad (n =0,1,2,\dots)}
\end{uncoverenv}
\EMgap
\begin{uncoverenv}<4->
\EMeq{X_n(x) =A_n\cos\lambda_nx,\qquad T_n(t) =C_ne^{-(c\lambda_n)^2t}}
\end{uncoverenv}
\end{EMbox}
\end{frame}

\def\EMtitlesize{26}\def\EMtitleskip{31}
\begin{frame}{حل معادلۀ دیفرانسیل انتقال حرارت یک‌بعدی (ادامه)}
\begin{EMbox}{ex}{1}{}
\begin{uncoverenv}<2->
\EMeq{u(x,t) =B_0+\sum_{n=1}^{\infty}B_ne^{-(c\lambda_n)^2t}\cos(\lambda_nx),\qquad \lambda_n =\frac{n\pi}L}
\end{uncoverenv}
\EMgap
\begin{uncoverenv}<3->
\EMp{(به ازای \EMim{\lambda_0=0} جملۀ ثابت \EMim{B_0} به دست می‌آید.) شرط اولیه \EMim{u(x,0)=B_0+\sum_{n=1}^{\infty}B_n\cos\!\left(\dfrac{n\pi}Lx\right)=f(x)}:}
\vspace{3bp}
\EMeq{B_0 =\frac1L\int_0^Lf(x)\,dx,\qquad B_n =\frac2L\int_0^Lf(x)\cos\!\left(\frac{n\pi}Lx\right)dx}
\end{uncoverenv}
\end{EMbox}
\end{frame}

\def\EMtitlesize{26}\def\EMtitleskip{31}
\begin{frame}{معادلۀ انتقال حرارت یک‌بعدی بر روی مفتول نامتناهی}
\begin{uncoverenv}<2->
\EMeq{\begin{cases}\dfrac{\partial u}{\partial t}=c^2\dfrac{\partial^2u}{\partial x^2}\\\noalign{\vskip 2mm} u(\infty,t)=u(-\infty,t)=0\\ u(x,0)=f(x),\quad-\infty<x<\infty\end{cases}}
\end{uncoverenv}
\EMgap
\begin{uncoverenv}<3->
\EMp{با فرض \EMim{u(x,t)=X(x)\times T(t)}:}
\vspace{3bp}
\EMeq{X(x)T'(t) =c^2X''(x)T(t)\  \Rightarrow\ \frac{T'}{c^2T} =\frac{X''}X =-\lambda^2\  \Rightarrow\ \begin{cases}X(x)=A_1(\lambda)\cos\lambda x+B_1(\lambda)\sin\lambda x=X_\lambda(x)\\ T(t)=C_1(\lambda)e^{-c^2\lambda^2t}=T_\lambda(t)\end{cases}}
\end{uncoverenv}
\EMgap
\begin{uncoverenv}<4->
\EMeq{u_\lambda(x,t) =\big[A(\lambda)\cos\lambda x+B(\lambda)\sin\lambda x\big]e^{-c^2\lambda^2t}\quad \Longrightarrow\quad u(x,t) \\[4pt] =\frac1\pi\int_0^\infty\big[A(\lambda)\cos\lambda x+B(\lambda)\sin\lambda x\big]e^{-c^2\lambda^2t}\,d\lambda}
\end{uncoverenv}
\end{frame}

\def\EMtitlesize{26}\def\EMtitleskip{31}
\begin{frame}{معادلۀ انتقال حرارت یک‌بعدی بر روی مفتول نامتناهی (ادامه)}
\begin{uncoverenv}<2->
\EMp{صحت معادله:}
\vspace{3bp}
\EMeq{\frac{\partial^2u}{\partial x^2} =\frac1\pi\int_0^\infty-\lambda^2\big[A\cos\lambda x+B\sin\lambda x\big]e^{-c^2\lambda^2t}\,d\lambda,\qquad \frac{\partial u}{\partial t} =\frac1\pi\int_0^\infty-c^2\lambda^2\big[A\cos\lambda x+B\sin\lambda x\big]e^{-c^2\lambda^2t}\,d\lambda\  \\[4pt] \Rightarrow\ \frac{\partial u}{\partial t} =c^2\frac{\partial^2u}{\partial x^2}\ \EMcheck}
\end{uncoverenv}
\EMgap
\begin{uncoverenv}<3->
\EMp{\textbf{اعمال شرط اولیه:} \EMim{u(x,0)=\dfrac1\pi\displaystyle\int_0^\infty\big[A(\lambda)\cos\lambda x+B(\lambda)\sin\lambda x\big]d\lambda=f(x)}، یعنی انتگرال فوریه (فصل ۳):}
\vspace{3bp}
\EMeq{A(\lambda) =\mathfrak{I}_C\{f(x)\} =\int_{-\infty}^{\infty}f(x)\cos\lambda x\,dx,\qquad B(\lambda) =\mathfrak{I}_S\{f(x)\} =\int_{-\infty}^{\infty}f(x)\sin\lambda x\,dx}
\end{uncoverenv}
\end{frame}

\def\EMtitlesize{26}\def\EMtitleskip{31}
\begin{frame}{معادلۀ انتقال حرارت یک‌بعدی بر روی مفتول نامتناهی (ادامه)}
\begin{uncoverenv}<2->
\EMeq{\begin{aligned}
u(x,t)&=\frac1\pi\int_0^\infty\left[\left(\int_{-\infty}^{\infty}f(\nu)\cos\lambda\nu\,d\nu\right)\cos\lambda x+\left(\int_{-\infty}^{\infty}f(\nu)\sin\lambda\nu\,d\nu\right)\sin\lambda x\right]e^{-c^2\lambda^2t}\,d\lambda\\
&=\frac1\pi\int_0^\infty\left\{\int_{-\infty}^{\infty}f(\nu)\big[\cos\lambda\nu\cos\lambda x+\sin\lambda\nu\sin\lambda x\big]d\nu\right\}e^{-c^2\lambda^2t}\,d\lambda\\
&=\frac1\pi\int_0^\infty\left\{\int_{-\infty}^{\infty}f(\nu)\cos\lambda(x-\nu)\,d\nu\right\}e^{-c^2\lambda^2t}\,d\lambda
\end{aligned}}
\end{uncoverenv}
\EMgap
\begin{EMbox}{imp}{2}{}
\begin{uncoverenv}<3->
\EMeq{u(x,t) =\frac1\pi\int_{-\infty}^{\infty}f(\nu)\left\{\int_0^\infty e^{-c^2\lambda^2t}\cos\lambda(x-\nu)\,d\lambda\right\}d\nu}
\end{uncoverenv}
\end{EMbox}
\end{frame}

\def\EMtitlesize{26}\def\EMtitleskip{31}
\begin{frame}{لم‌های ریاضی}
\begin{EMbox}{thm}{1}{}
\begin{uncoverenv}<2->
\EMeq{I =\int_0^\infty e^{-x^2}\,dx =\frac{\sqrt\pi}2}
\end{uncoverenv}
\end{EMbox}
\begin{EMbox}{prf}{2}{}
\begin{uncoverenv}<3->
\EMeq{I^2 =\int_0^\infty e^{-x^2}dx\times\int_0^\infty e^{-y^2}dy =\int_0^\infty\!\!\int_0^\infty e^{-(x^2+y^2)}\,dx\,dy\ \overset{\rho^2=x^2+y^2}{=}\ \int_0^{\pi/2}\!\!\int_0^\infty\rho\,e^{-\rho^2}\,d\rho\,d\theta \\[4pt] =\int_0^{\pi/2}\left[\frac{-1}2e^{-\rho^2}\right]_0^\infty d\theta =\frac\pi4}
\end{uncoverenv}
\end{EMbox}
\begin{EMbox}{thm}{3}{}
\begin{uncoverenv}<4->
\EMeq{\int_0^\infty e^{-s^2}\cos(2bs)\,ds =\frac{\sqrt\pi}2e^{-b^2}}
\end{uncoverenv}
\end{EMbox}
\begin{EMbox}{prf}{4}{}
\begin{uncoverenv}<5->
\EMeq{I(b) =\int_0^\infty e^{-s^2}\cos(2bs)\,ds\  \Rightarrow\ \frac{dI(b)}{db} =-\int_0^\infty2se^{-s^2}\sin(2bs)\,ds}
\end{uncoverenv}
\end{EMbox}
\end{frame}

\def\EMtitlesize{26}\def\EMtitleskip{31}
\begin{frame}{لم‌های ریاضی (ادامه)}
\begin{EMbox}{prf}{1}{}
\begin{uncoverenv}<2->
\EMp{با انتگرال‌گیری جزء به جزء (\EMim{du=-2se^{-s^2}ds}، \EMim{v=\sin2bs}):}
\vspace{3bp}
\EMeq{\frac{dI(b)}{db} =e^{-s^2}\sin2bs\Big|_0^\infty-2b\int_0^\infty e^{-s^2}\cos2bs\,ds =-2b\,I(b)\  \Rightarrow\ I(b) =ke^{-b^2},\qquad I(0) =k \\[4pt] =\int_0^\infty e^{-s^2}ds =\frac{\sqrt\pi}2}
\end{uncoverenv}
\end{EMbox}
\begin{uncoverenv}<3->
\EMp{با تغییر متغیر \EMim{c\lambda\sqrt t=s}:}
\vspace{3bp}
\EMeq{\int_0^\infty e^{-c^2\lambda^2t}\cos\lambda(x-\nu)\,d\lambda =\frac1{c\sqrt t}\int_0^\infty e^{-s^2}\cos\!\left[2\frac{(x-\nu)}{2c\sqrt t}s\right]ds =\frac1{c\sqrt t}\frac{\sqrt\pi}2e^{-\left(\frac{x-\nu}{2c\sqrt t}\right)^2}}
\end{uncoverenv}
\EMgap
\begin{EMbox}{imp}{3}{}
\begin{uncoverenv}<4->
\EMeq{u(x,t) =\frac1{2c\sqrt{\pi t}}\int_{-\infty}^{\infty}f(\nu)\,e^{-\frac{(x-\nu)^2}{4c^2t}}\,d\nu}
\end{uncoverenv}
\end{EMbox}
\begin{uncoverenv}<5->
\EMp{با تغییر متغیر \EMim{z=\dfrac{\nu-x}{2c\sqrt t}} (یعنی \EMim{\nu=x+2c\sqrt t\,z} و \EMim{d\nu=2c\sqrt t\,dz}):}
\end{uncoverenv}
\end{frame}

\def\EMtitlesize{26}\def\EMtitleskip{31}
\begin{frame}{لم‌های ریاضی (ادامه)}
\begin{EMbox}{imp}{1}{}
\begin{uncoverenv}<2->
\EMeq{u(x,t) =\frac1{\sqrt\pi}\int_{-\infty}^{\infty}f\big(x+2c\sqrt t\,z\big)\,e^{-z^2}\,dz}
\end{uncoverenv}
\end{EMbox}
\end{frame}

\def\EMtitlesize{26}\def\EMtitleskip{31}
\begin{frame}{تابع خطا}
\begin{minipage}[t]{500bp}\vspace{0pt}
\begin{uncoverenv}<2->
\EMp{در ریاضی اغلب تابع خطا را به صورت زیر تعریف می‌کنند:}
\vspace{3bp}
\EMeq{\operatorname{erf}(x)\triangleq\frac2{\sqrt\pi}\int_0^xe^{-t^2}\,dt}
\end{uncoverenv}
\EMgap
\begin{uncoverenv}<3->
\EMp{خواص تابع خطا:}
\end{uncoverenv}
\EMgap
\begin{uncoverenv}<4->
\EMitem{\lr{\textbullet}}{\dimexpr\linewidth-26bp\relax}{خاصیت تقارن فرد: \EMim{\operatorname{erf}(-x)=-\operatorname{erf}(x)}}
\end{uncoverenv}
\EMgap
\begin{uncoverenv}<5->
\EMitem{\lr{\textbullet}}{\dimexpr\linewidth-26bp\relax}{خاصیت مجانبی: \EMim{\operatorname{erf}(\infty)=1}}
\end{uncoverenv}
\EMgap
\begin{uncoverenv}<6->
\EMitem{\lr{\textbullet}}{\dimexpr\linewidth-26bp\relax}{بسط مک‌لورن این تابع به صورت زیر است (با بسط \EMim{e^x} شروع و این تساوی را ثابت کنید):}
\vspace{3bp}
\EMeq{\operatorname{erf}(x) =\frac2{\sqrt\pi}\left[x-\frac{x^3}{1!\times3}+\frac{x^5}{2!\times5}-\frac{x^7}{3!\times7}+-\cdots\right] \\[4pt] =\frac2{\sqrt\pi}\sum_{n=1}^{\infty}(-1)^{n+1}\frac{x^{2n-1}}{(n-1)!\,(2n-1)}}
\end{uncoverenv}
\EMgap
\end{minipage}
\EMside{221.6bp}{\uncover<7->{\EMfigure{1.389}{m05-erf}{نمودار تابع خطا \EMim{\operatorname{erf}(x)} برای \EMim{-2\le x\le2}}}}
\end{frame}

\def\EMtitlesize{26}\def\EMtitleskip{31}
\begin{frame}{تابع خطا (ادامه)}
\begin{EMbox}{ex}{1}{}
\EMp{جواب معادلۀ بالا را با فرض شرط اولیۀ مقابل به دست آورید.}
\vspace{3bp}
\EMeq{f(x) =\begin{cases}U_0&|x|<1\\ 0&|x|\ge1\end{cases}}
\EMgap
\begin{uncoverenv}<2->
\EMp{از شکل دوم جواب استفاده می‌کنیم:}
\vspace{3bp}
\EMeq{u(x,t) =\frac{U_0}{2c\sqrt{\pi t}}\int_{-1}^{1}e^{-\frac{(x-\nu)^2}{4c^2t}}\,d\nu\ \overset{z=\frac{\nu-x}{2c\sqrt t}}{=}\ \frac{U_0}{\sqrt\pi}\int_{-\frac{1+x}{2c\sqrt t}}^{\frac{1-x}{2c\sqrt t}}e^{-z^2}\,dz}
\end{uncoverenv}
\EMgap
\begin{uncoverenv}<3->
\EMeq{=\frac{U_0}{\sqrt\pi}\left\{\int_{-\frac{1+x}{2c\sqrt t}}^{0}e^{-z^2}dz+\int_0^{\frac{1-x}{2c\sqrt t}}e^{-z^2}dz\right\} =\frac{U_0}{\sqrt\pi}\frac{\sqrt\pi}2\left\{\frac2{\sqrt\pi}\int_0^{\frac{1+x}{2c\sqrt t}}e^{-z^2}dz+\frac2{\sqrt\pi}\int_0^{\frac{1-x}{2c\sqrt t}}e^{-z^2}dz\right\}}
\end{uncoverenv}
\end{EMbox}
\end{frame}

\def\EMtitlesize{26}\def\EMtitleskip{31}
\begin{frame}{تابع خطا (ادامه)}
\begin{minipage}[t]{500bp}\vspace{0pt}
\begin{EMbox}{ex}{1}{}
\begin{uncoverenv}<2->
\EMeq{u(x,t) =\frac{U_0}2\left[\operatorname{erf}\!\left(\frac{1+x}{2c\sqrt t}\right)+\operatorname{erf}\!\left(\frac{1-x}{2c\sqrt t}\right)\right]}
\end{uncoverenv}
\end{EMbox}
\end{minipage}
\EMside{82.5bp}{\uncover<3->{\EMfigure{1.193}{m05-heat-pulse}{توزیع دما \EMim{u(x,t)} در لحظه‌های \EMim{t=0}، \EMim{\frac18}، \EMim{\frac12}، \EMim{1}، \EMim{2} و \EMim{8} برای پالس اولیۀ \EMim{U_0=100} و \EMim{c=1}}}}
\end{frame}

\def\EMtitlesize{26}\def\EMtitleskip{31}
\begin{frame}{تابع خطا (ادامه)}
\begin{EMbox}{ex}{1}{}
\EMp{جواب معادلۀ بالا را با فرض شرط اولیۀ مقابل به دست آورید.}
\vspace{3bp}
\EMeq{f(x) =\begin{cases}1&x>0\\ 0&x\le0\end{cases}}
\EMgap
\begin{uncoverenv}<2->
\EMp{از شکل دوم جواب استفاده می‌کنیم:}
\vspace{3bp}
\EMeq{u(x,t) =\frac1{2c\sqrt{\pi t}}\int_0^\infty e^{-\frac{(x-\nu)^2}{4c^2t}}\,d\nu\ \overset{z=\frac{\nu-x}{2c\sqrt t}}{=}\ \frac1{\sqrt\pi}\int_{-\frac x{2c\sqrt t}}^{\infty}e^{-z^2}\,dz =\frac1{\sqrt\pi}\left\{\int_{-\frac x{2c\sqrt t}}^{0}e^{-z^2}dz+\int_0^\infty e^{-z^2}dz\right\}}
\end{uncoverenv}
\EMgap
\begin{uncoverenv}<3->
\EMeq{u(x,t) =\frac12\left[\operatorname{erf}\!\left(\frac x{2c\sqrt t}\right)+1\right]}
\end{uncoverenv}
\end{EMbox}
\end{frame}
